\chapter{Introduction and Problem}\label{ch:intro}

\statcard{\$685\,B--\$890\,B}{Returns of the US retail industry in 2024, by two different methods}{Sources: Appriss Retail \cite{appriss2024}; NRF \cite{nrf2024}}
\statcard{24.4\%}{Online apparel orders returned in the 12 months to March 2023 (US survey of 100 decision-makers)}{Source: Coresight Research \cite{coresight2023}}
\statcard{53\%}{Surveyed decision-makers who name size or fit the top reason for apparel returns}{Source: Coresight Research \cite{coresight2023}; share of respondents, not of returns}
\statcard{17\%}{Cost of one return as a share of the prime cost (one Scandinavian footwear retailer)}{Source: Gustafsson et al. \cite{gustafsson2021}}

\section{The returns bill}\label{sec:returns-bill}

Every order that comes back is a sale that did not finish and a cost that did. The two best estimates of what this costs retailers in the United States are large and differ from each other. Appriss Retail reports that total returns of the retail industry amounted to \$685 billion in merchandise in 2024, and the National Retail Federation projects \$890 billion for the same year \cite{appriss2024,nrf2024}. The gap of \$205 billion is a difference of method and not a contradiction. The Appriss figure is built from surveys of 150 North American retail executives and 1,000 consumers together with data of the US Census Bureau, while the NRF figure is a projection published on the federation's report page \cite{appriss2024,nrf2024}. We therefore quote the range and not a single number, and we do not add the two. Appriss also reports that online sales account for \$362 billion of the returned value and that fraud and abuse account for 15.14\% of all returns \cite{appriss2024}. The second figure matters for honesty: not every return is a fit problem, and a fitting tool cannot reduce fraud.

Two different return rates circulate in the trade press, and they must not be merged. Appriss reports that online returns reached 24.52\% of all online sales, which covers every product category \cite{appriss2024}. Coresight Research estimates that the average return rate of online apparel orders was 24.4\% for the twelve months ended March 2023, from a survey of 100 decision-makers at US apparel brands and retailers \cite{coresight2023}. The similarity of the two numbers is a coincidence of two scopes. This plan uses the Coresight figure, because it is the one that is specific to apparel, and it states each time that the base is a small US survey.

\section{Why fit is the core cause}\label{sec:fit-cause}

Among the reasons for apparel returns, size and fit come first. In the Coresight survey, 53\% of respondents name size or fit as the top reason for online apparel returns, followed by color (16\%) and damage (10\%) \cite{coresight2023}. This is a share of respondents and not a share of returns, and the plan treats it that way: our own model assumes that 40\% of returns are fit-related, which is a deliberate step below the survey figure (Chapter~\ref{ch:finance}). The same picture appears in the research literature. Abdulla et al. state that online shopping yields a large percentage of returns due to size and fit \cite{abdulla2019}, and Alemany et al. describe the high number of returns and their management, logistics and operating costs as an important economic burden for fashion companies \cite{alemany2013}. Peng and Duraffourg report that the return rate in the textile industry is in general between 20\% and 40\%, a secondary claim that they make without a citation in their abstract \cite{peng2023}.

The literature also tells us what kind of problem this is. A shopper who orders online cannot try the garment on, so the decision is made on incomplete information: the shopper knows his or her body only roughly, and the garment is described by a size label that means different things at different brands. Alemany et al. note that earlier fitting technologies did not reach a compromise between cost, efficacy and accuracy that would allow wide market penetration \cite{alemany2013}, and we treat this as the historical gap that a lower-cost scan and a data-driven model can address. In a different field, Pfaff and Spann report that augmented reality is particularly effective for larger products with higher spatial fit uncertainty \cite{pfaff2025}. We cite it only as an analogy: the study is about home interiors and not about apparel.

\section{What one return costs}\label{sec:return-cost}

A return is expensive because it is handled twice. In the one study that built the cost of a return from transaction data, 2,229 return transactions of a Scandinavian fashion footwear retailer, the cost of a return is approximately 17\% of the prime cost, and product handling and transportation together make up 72\% of that cost \cite{gustafsson2021}. This is a single retailer and a single product category, so it is an anchor for the model and not a market constant. Coresight estimates, for the United States, that a 24.4\% return rate on a \$155.8 billion online apparel and footwear market translates to \$38 billion in returns and \$25.1 billion in processing costs in 2023 \cite{coresight2023}. We use the Gustafsson ratio in our model and show the Coresight figures only as a statement of scale.

For Israel, which is the first market of this plan, the legal setting makes returns structurally material. According to the official guide of the Israeli Consumer Protection and Fair Trade Authority, a consumer may cancel a distance-selling transaction within fourteen days of closing the transaction or of receiving the document with its details, whichever is later; use of the product does not remove the right, and the business must refund the consumer within fourteen days of receiving the cancellation notice \cite{il_cpfta_guide}. The business may charge a cancellation fee of 5\% of the transaction or NIS~100, whichever is lower, and may not charge it when the product is defective or differs from the details provided \cite{il_cpfta_guide}. We found no verified Israeli return rate. The plan's return rate therefore stays at the US figure of \assumed{24.4\%}, and it is an assumption anchored on a US survey and not Israeli data.

\section{The thesis and how this plan tests it}\label{sec:thesis}

ClosetTwin is a business-to-business virtual fitting room for fashion e-commerce. The shopper gives a body scan from a phone, the retailer supplies the garment's measurements, and the service combines both with the shopper's stated fit preference to recommend a size. The claim at the center of the plan is a hypothesis and not a fact: \emph{a fit tool built on scan, garment data and preference lowers fit-related returns enough to pay for itself}. The evidence is mixed. Gustafsson et al. model reductions of fit-related return costs of up to 80\% under assumed verification rates, but this is a scenario \cite{gustafsson2021}. The largest field study in our sources, covering 496,365 items ordered by 75,707 customers, found that customers who used a size finder were 0.65\% more likely to return an item, although their customer lifetime value was higher \cite{patel2025}. A plan that ignored the second result would be dishonest.

The plan therefore treats the effect as the quantity to be measured. Chapter~\ref{ch:market} sizes the market and finds that Israel alone is too small to carry the company. Chapter~\ref{ch:competition} places ClosetTwin among the existing vendors. Chapters~\ref{ch:product} and~\ref{ch:business} describe the product and its pricing, which is set at a fraction of the value a retailer saves if the hypothesis holds. Chapter~\ref{ch:gtm} designs a randomized pilot that can accept or reject the hypothesis, and Chapters~\ref{ch:finance} to~\ref{ch:conclusion} show what the company costs, what it earns and what it asks for if the pilot passes. Where a number is our own, it is marked as derived or assumed, and the legend is \ctlegend{}.

\ctchapterend
